\documentclass[12pt,a4paper]{article}

\usepackage[utf8]{inputenc}
\usepackage[T1]{fontenc}
\usepackage[brazil]{babel}
\usepackage{lmodern}
\usepackage{geometry}
\usepackage{microtype}
\usepackage{array}
\usepackage{tabularx}
\usepackage{enumitem}
\usepackage{xcolor}
\usepackage{fancyhdr}

\geometry{top=1.7cm,bottom=1.7cm,left=1.8cm,right=1.8cm}
\setlength{\parindent}{0pt}
\setlength{\parskip}{0.45em}
\setlist[itemize]{leftmargin=*,itemsep=0.25em,topsep=0.2em}
\definecolor{azul}{HTML}{173B57}
\definecolor{cinza}{HTML}{F2F4F7}

\pagestyle{fancy}
\setlength{\headheight}{15pt}
\fancyhf{}
\lhead{Sequência didática de \textit{machine learning}}
\rhead{Aula 3: regressão logística}
\cfoot{\thepage}
\renewcommand{\headrulewidth}{0.3pt}

\newcommand{\linha}{\rule{\linewidth}{0.35pt}}
\newcommand{\bloco}[1]{%
  \vspace{0.25em}
  \noindent\colorbox{cinza}{\parbox{0.96\linewidth}{#1}}%
  \vspace{0.3em}
}

\begin{document}

\begin{center}
  {\LARGE\bfseries\color{azul} Aula 3: regressão logística}\par
  \vspace{0.25em}
  {\large Probabilidade, limiar e classificação}\par
  \vspace{0.35em}
  {\small Roteiro de aplicação da sequência didática}\par
\end{center}

\bloco{\textbf{Lugar desta aula na pesquisa.} Este é o terceiro encontro da sequência de oito aulas. A atividade trata da diferença entre uma probabilidade calculada e a categoria exibida após a aplicação de uma regra. A exploração individual tem função formativa; o professor registra dificuldades recorrentes da turma.}

\section*{1. Identificação}

\begin{tabularx}{\linewidth}{@{}>{\bfseries}p{4.1cm}X@{}}
Turma: & \rule{8.5cm}{0.35pt} \\
Data: & \rule{8.5cm}{0.35pt} \\
Duração: & 50 minutos \\
Organização: & Exploração individual e síntese oral coletiva. \\
Recursos: & Plataforma no módulo de regressão logística, computadores do laboratório ou Chromebooks, quadro e folha de atividade. \\
\end{tabularx}

\section*{2. Objetivos da aula}

Ao final do encontro, espera-se que os estudantes distingam probabilidade de categoria final, identifiquem o limiar como regra de decisão e expliquem que uma mudança no limiar pode alterar a classificação sem modificar a probabilidade calculada.

\textbf{Relação curricular:} EM13CO10, EM13CO11 e EM13CO12.

\textbf{Conceitos em foco:} dado, probabilidade, classificação binária, limiar, regra de decisão, categoria e limite de previsão.

\section*{3. Preparação antes da aula}

\begin{itemize}
  \item Abrir o módulo de regressão logística e confirmar o cenário sintético de alerta de irrigação.
  \item Conferir se a probabilidade calculada, o limiar e a categoria final aparecem de forma legível.
  \item Confirmar que cada estudante terá acesso individual ao módulo, em computador do laboratório ou Chromebook.
  \item Reservar um espaço no quadro para comparar a mesma probabilidade sob limiares diferentes.
\end{itemize}

\bloco{\textbf{Cuidado didático.} O alerta de irrigação é uma situação sintética. A probabilidade não deve ser apresentada como certeza sobre acontecimentos reais, pois ela depende dos dados disponíveis e da forma como o modelo foi construído.}

\section*{4. Desenvolvimento da aula}

\renewcommand{\arraystretch}{1.22}
\begin{tabularx}{\linewidth}{|>{\centering\arraybackslash}p{1.65cm}|p{5.0cm}|X|}
\hline
\textbf{Tempo} & \textbf{Ação com a turma} & \textbf{Mediação docente} \\
\hline
0--7 min & Retomar a regressão linear e contrastar previsão numérica com decisão entre categorias. & Perguntar: ``Na aula anterior, a saída era um número. Que tipo de saída teremos quando o sistema precisa mostrar um alerta ou não mostrar?'' \\
\hline
7--12 min & Apresentar o cenário fictício de uma estufa. Mostrar uma probabilidade de irrigação e a regra inicial para exibir o alerta. & Confirmar que a porcentagem representa uma chance calculada e que o alerta é uma decisão tomada a partir de uma regra. \\
\hline
12--18 min & Pedir que cada estudante antecipe o efeito de mudar somente o limiar de 60\% para 80\%, mantendo a probabilidade em 70\%. & Solicitar que indiquem separadamente o que pode mudar: a porcentagem, a categoria final ou ambos. \\
\hline
18--33 min & Explorar o módulo de forma guiada. Alterar o limiar e observar a categoria; comparar casos próximos da regra. & Distinguir a curva ou probabilidade calculada da decisão produzida pelo limiar. Não tratar a classificação como verdade sobre a estufa. \\
\hline
33--43 min & Retomar previsões individuais e comparar 60\% e 80\% para a mesma probabilidade. & Pedir uma explicação completa: valor calculado, regra usada e categoria exibida. Registrar no quadro exemplos em que a categoria muda. \\
\hline
43--50 min & Realizar a atividade individual breve e encerrar com uma síntese oral. & Esclarecer o enunciado sem indicar a resposta. Retomar que modelos dependem de dados, regras e escolhas humanas. \\
\hline
\end{tabularx}

\section*{5. Perguntas para conduzir a conversa}

\begin{itemize}
  \item ``O que significa uma chance calculada de 70\%?''
  \item ``Qual regra define quando o alerta aparece?''
  \item ``O que muda quando alteramos apenas o limiar?''
  \item ``A probabilidade calculada mudou junto com a categoria?''
  \item ``Por que uma classificação não deve ser tratada como certeza?''
\end{itemize}

\section*{6. Evidência a observar}

Observar se os estudantes separam a probabilidade calculada da decisão final e se conseguem explicar a função do limiar. No registro docente, anotar confusões entre porcentagem e categoria, justificativas formuladas individualmente e condições de uso da plataforma que tenham afetado a aula.

\newpage

\begin{center}
  {\LARGE\bfseries\color{azul} Atividade individual breve}\par
  \vspace{0.2em}
  {\large Aula 3: probabilidade, limiar e alerta}\par
  \vspace{0.35em}
  {\small Material formativo para o final da aula}\par
\end{center}

\bloco{\textbf{Orientação ao estudante.} Esta atividade não vale nota. Responda individualmente e explique como você pensou. Use o que foi observado na plataforma e o que foi discutido na aula.}

\textbf{Código do participante:} \rule{5cm}{0.4pt}
\hfill
\textbf{Data:} \rule{3cm}{0.4pt}

\section*{Situação}

Um sistema fictício calcula a chance de uma estufa precisar de irrigação. Em uma leitura, a chance calculada foi de 70\%. O sistema mostra o alerta \textit{irrigar} somente quando a chance é igual ou maior que o limiar definido.

\section*{Questão 1}

Com limiar de 60\%, qual alerta será mostrado? E com limiar de 80\%? Explique por que as respostas podem ser diferentes mesmo com a mesma chance calculada.

\vspace{0.25em}
\linha\par\vspace{0.75em}
\linha\par\vspace{0.75em}
\linha

\section*{Questão 2}

Ao mudar somente o limiar de 60\% para 80\%, a chance calculada de 70\% também muda? Explique a diferença entre chance calculada e alerta mostrado.

\vspace{0.25em}
\linha\par\vspace{0.75em}
\linha\par\vspace{0.75em}
\linha

\section*{Questão 3}

Uma chance de 70\% significa que a estufa certamente precisa de irrigação? Explique sua resposta e indique uma informação que poderia ajudar a interpretar melhor esse alerta.

\vspace{0.25em}
\linha\par\vspace{0.75em}
\linha

\end{document}
